\section{Background and Motivation}

\subsection{Asymmetric Quantization for KV Cache}

\kivi{}~\cite{kivi} adopts asymmetric (affine) quantization where each value maps as $x = s(\hat{x} - z)$ with per-channel scale $s$ and zero-point $z$.
For 2-bit quantization, $\hat{x} \in \{0,1,2,3\}$ requiring only 0.25 bytes per element --- an $8\times$ reduction from float16.

The attention score between query $Q$ and key $K_i$ with asymmetric quantization expands as:
\begin{equation}
Q \cdot K_i = s_Q s_K \!\left[\sum_{j=1}^{d} \hat{q}_j \hat{k}_{ij} - z_K\!\sum_{j=1}^{d} \hat{q}_j - z_Q\!\sum_{j=1}^{d} \hat{k}_{ij} + d z_Q z_K\right]
\label{eq:full}
\end{equation}
This reveals four operations: (1) main dot product $\sum \hat{q}\hat{k}$ on 2-bit data, (2) Q-reduction $\sum\hat{q}$, (3) K-reduction $\sum\hat{k}$, (4) constant $d z_Q z_K$.

\textbf{Key observation:} operations (1)--(3) operate on quantized 2-bit data, processable by simple near-memory compute units. Only the final scaling $s_Q s_K$ requires float16.

\subsection{Near-Bank PIM Architecture}

Near-bank PIM places lightweight PEs inside each DRAM bank~\cite{seshadri2017ambit}.
Each bank has: a row buffer ($R$ bytes) serving as PE working memory, a MAC unit processing $W$ bytes per cycle at frequency $f$, and local DRAM-to-rowbuffer bandwidth $B$ GB/s.

For HBM3E with 8 stacks $\times$ 256 PIM banks, $N=2048$ banks operate in parallel.
Data is striped across banks; each bank processes its partition independently.

\subsection{Limitations of Existing Approaches}

\textbf{GPU-only:} \kivi{}'s dequantization and reduction overhead on GPU limits throughput as extra terms in Eq.~\ref{eq:full} compete for compute units.

\textbf{Logic-die PIM~\cite{duplex}:} Compute on separate logic die requires data traversal via TSVs, introducing bandwidth bottlenecks for fine-grained bank-local operations.

\textbf{PIM dequantization:} Dequantizing 2-bit to float16 inside PIM ($s \cdot (\hat{x} - z)$) adds $3\times$ PE time per element, doubling Q@K latency. This motivates our approach of \emph{decomposing and keeping operations in the 2-bit domain}.
